\documentclass[a4paper,10pt]{article}
\usepackage{lamia-pedagogy}

\begin{document}

% =========================================================================
% PAGE 1 : Fiche de Cours - Le Présent du 1er Groupe
% =========================================================================
\ficheHeader{1 Heure}{\makebox[2.5cm]{\dotfill}}{1\textsuperscript{ère} AM}{2026 / 2027}

\vspace{-3pt}
\competenceBanner{\faBookOpen\ Projet 01 / Séquence 01 — Séance 06 : Conjugaison}

\vspace{-2pt}
\modaliteSupportBar{Observation collective puis entraînement individuel}{Support : Phrases d'hygiène / Cahier de conjugaison}

\vspace{1pt}

% Phrases supports
\begin{tcolorbox}[
  enhanced,
  colback=primarylight!40,
  colframe=primary,
  arc=2mm,boxrule=0.8pt,
  top=4pt,bottom=4pt,left=8pt,right=8pt,
  title={\small\textbf{\textcolor{white}{\faEye\ Phrases d'observation}}},
  coltitle=white,colbacktitle=primary
]
{\small
\begin{itemize}[leftmargin=*,noitemsep,topsep=1pt]
  \item Maman \textbf{lave} le linge avec une lessive parfumée.
  \item Tu \textbf{brosses} tes dents après chaque repas.
  \item Nous \textbf{nettoyons} la maison pour éliminer la poussière.
  \item Les médecins \textbf{conseillent} de se laver les mains régulièrement.
\end{itemize}
}
\end{tcolorbox}

\vspace{3pt}

% Table de conjugaison
\begin{minipage}[t]{0.52\textwidth}
\begin{tcolorbox}[
  enhanced,
  colback=secondarylight!50,
  colframe=secondary,
  arc=2mm,boxrule=0.8pt,
  top=4pt,bottom=4pt,left=7pt,right=7pt,
  title={\small\textbf{\textcolor{secondary!90!black}{\faStar\ Règle de Conjugaison (-er)}}},
  colbacktitle=white,coltitle=secondary!90!black,
  attach boxed title to top left={yshift=-2mm,xshift=4mm},
  boxed title style={boxrule=0.6pt,colframe=secondary,arc=1.2mm,top=1pt,bottom=1pt}
]
\vspace{2pt}
{\small
Pour conjuguer un verbe du 1\textsuperscript{er} groupe au présent :\\
On prend le \textbf{radical} et on ajoute les terminaisons :

\vspace{3pt}
\begin{tabularx}{\linewidth}{l X}
\textbf{je} & radical $+$ \textbf{\textcolor{secondary}{\large -e}} \\
\textbf{tu} & radical $+$ \textbf{\textcolor{secondary}{\large -es}} \\
\textbf{il / elle / on} & radical $+$ \textbf{\textcolor{secondary}{\large -e}} \\
\textbf{nous} & radical $+$ \textbf{\textcolor{secondary}{\large -ons}} \\
\textbf{vous} & radical $+$ \textbf{\textcolor{secondary}{\large -ez}} \\
\textbf{ils / elles} & radical $+$ \textbf{\textcolor{secondary}{\large -ent}} \\
\end{tabularx}
}
\end{tcolorbox}
\end{minipage}%
\hfill
\begin{minipage}[t]{0.45\textwidth}
\begin{tcolorbox}[
  enhanced,
  colback=accentlight!35,
  colframe=accent,
  arc=2mm,boxrule=0.8pt,
  top=4pt,bottom=4pt,left=7pt,right=7pt,
  title={\small\textbf{\textcolor{accent!90!black}{\faLightbulb\ Modèle : Laver}}},
  colbacktitle=white,coltitle=accent!90!black,
  attach boxed title to top left={yshift=-2mm,xshift=4mm},
  boxed title style={boxrule=0.6pt,colframe=accent,arc=1.2mm,top=1pt,bottom=1pt}
]
\vspace{2pt}
{\small
\textbf{Verbe LAVER au présent :}
\begin{itemize}[leftmargin=*,noitemsep,topsep=1pt]
  \item je lav\textbf{\textcolor{accent!90!black}{e}}
  \item tu lav\textbf{\textcolor{accent!90!black}{es}}
  \item il, elle lav\textbf{\textcolor{accent!90!black}{e}}
  \item nous lav\textbf{\textcolor{accent!90!black}{ons}}
  \item vous lav\textbf{\textcolor{accent!90!black}{ez}}
  \item ils, elles lav\textbf{\textcolor{accent!90!black}{ent}}
\end{itemize}
\par\noindent{\tiny\textit{Attention : Verbes en -ger (nous mangeons), en -cer (nous commençons).}}
}
\end{tcolorbox}
\end{minipage}

\vspace{3pt}

% Synthèse
\begin{tcolorbox}[
  enhanced,
  colback=white,
  colframe=cardborder,
  arc=1.5mm,boxrule=0.6pt,
  top=3pt,bottom=3pt,left=7pt,right=7pt
]
{\footnotesize
\faExclamationCircle\ \textbf{Règle d'or :} Les verbes du 1\textsuperscript{er} groupe se terminent tous par \textbf{-er} à l'infinitif (sauf le verbe \textit{aller} qui est du 3\textsuperscript{ème} groupe). Leurs terminaisons au présent sont régulières et invariables.
}
\end{tcolorbox}

\newpage

% =========================================================================
% PAGE 2 : Fiche d'Exercices d'Application
% =========================================================================
\ficheHeader{1 Heure}{Nom : \makebox[4cm]{\dotfill}}{1\textsuperscript{ère} AM}{Classe : \makebox[1.5cm]{\dotfill}}

\vspace{-3pt}
\competenceBanner{\faPenFancy\ Fiche d'Exercices : Le Présent des Verbes en -ER}

\vspace{2pt}

% ACTIVITÉ 1 : Terminaisons & Infinitif
\begin{tcolorbox}[
  enhanced,
  colback=white,
  colframe=cardborder,
  arc=2mm,boxrule=0.8pt,
  top=4pt,bottom=4pt,left=8pt,right=8pt,
  title={\small\textbf{\textcolor{primary}{\faPen\ Activité 01 : Ajoute la bonne terminaison et donne l'infinitif du verbe}} \hfill \badge[primary]{Terminaisons}},
  coltitle=primary,colbacktitle=primarylight,titlerule=0.5pt,titlerule style=cardborder
]
{\small
\begin{enumerate}[label=\textbf{\arabic*.},leftmargin=*,noitemsep,topsep=2pt]
  \item Tu bross\drule{1cm} tes dents trois fois par jour. $\longrightarrow$ Infinitif : \drule{3cm}
  \item L'infirmière nettoy\drule{1cm} la plaie avec du désinfectant. $\longrightarrow$ Infinitif : \drule{3cm}
  \item Nous utilis\drule{1.2cm} du savon doux pour le visage. $\longrightarrow$ Infinitif : \drule{3cm}
  \item Les enfants rinç\drule{1.2cm} leurs mains sous l'eau tiède. $\longrightarrow$ Infinitif : \drule{3cm}
\end{enumerate}
}
\end{tcolorbox}

\vspace{2pt}

% ACTIVITÉ 2 : Conjugaison
\begin{tcolorbox}[
  enhanced,
  colback=white,
  colframe=cardborder,
  arc=2mm,boxrule=0.8pt,
  top=4pt,bottom=4pt,left=8pt,right=8pt,
  title={\small\textbf{\textcolor{primary}{\faEdit\ Activité 02 : Mets les verbes entre parenthèses au présent de l'indicatif}} \hfill \badge[secondary]{Conjugaison}},
  coltitle=primary,colbacktitle=primarylight,titlerule=0.5pt,titlerule style=cardborder
]
{\small
\begin{enumerate}[label=\textbf{\arabic*.},leftmargin=*,noitemsep,topsep=2pt]
  \item Je (\textbf{frotter}) \drule{3cm} mes paumes avec insistance.
  \item Vous (\textbf{sécher}) \drule{3cm} vos mains avec une serviette propre.
  \item Nous (\textbf{éliminer}) \drule{3cm} la saleté grâce au savon.
  \item Ils (\textbf{protéger}) \drule{3cm} leur corps contre les maladies.
\end{enumerate}
}
\end{tcolorbox}

\vspace{2pt}

% ACTIVITÉ 3 : Transformation du sujet
\begin{tcolorbox}[
  enhanced,
  colback=white,
  colframe=cardborder,
  arc=2mm,boxrule=0.8pt,
  top=4pt,bottom=4pt,left=8pt,right=8pt,
  title={\small\textbf{\textcolor{primary}{\faSync*\ Activité 03 : Réécris chaque phrase avec le nouveau sujet}} \hfill \badge[accent]{Accord Sujet-Verbe}},
  coltitle=primary,colbacktitle=primarylight,titlerule=0.5pt,titlerule style=cardborder
]
{\small
\begin{enumerate}[label=\textbf{\arabic*.},leftmargin=*,noitemsep,topsep=2pt]
  \item L'élève \textbf{change} sa brosse à dents régulièrement.\\
    $\longrightarrow$ \textbf{Nous} \drule{11cm}
  \item Je \textbf{respecte} les règles d'hygiène corporelle.\\
    $\longrightarrow$ \textbf{Les enfants} \drule{10cm}
\end{enumerate}
}
\end{tcolorbox}

\vspace{2pt}

% ACTIVITÉ 4 : Production de phrases
\begin{tcolorbox}[
  enhanced,
  colback=white,
  colframe=cardborder,
  arc=2mm,boxrule=0.8pt,
  top=4pt,bottom=4pt,left=8pt,right=8pt,
  title={\small\textbf{\textcolor{primary}{\faPencilRuler\ Activité 04 : Rédige deux phrases à l'aide des verbes donnés}} \hfill \badge[softred]{Production}},
  coltitle=primary,colbacktitle=primarylight,titlerule=0.5pt,titlerule style=cardborder
]
{\small
\textbf{Consigne :} Conjugue au présent les verbes \textbf{laver} et \textbf{brosser} pour parler de ton hygiène :\\
1. (\textbf{laver}) : \drule{14.5cm}\\[8pt]
2. (\textbf{brosser}) : \drule{14.5cm}
}
\end{tcolorbox}

\end{document}
